% 第2章 自旋交换
\chapter{自旋交换}

当固体中许多电子的磁矩排成一致时，便出现铁磁性。反铁磁性和自旋密度波描述磁矩的振荡式排列。电子磁矩之间的经典偶极相互作用（量级 $10^{-5}$~eV）远远不足以解释观测到的磁转变温度（过渡金属和稀土化合物中量级为 $10^{2}\text{--}10^{3}$~K）。因此，量子力学发展初期人们就认识到，产生磁性的耦合机制来源于电子的如下基本性质：

\begin{enumerate}[itemsep=1pt]
\item 电子的自旋；
\item 电子的动能（离域化能量）；
\item Pauli 不相容原理（Fermi 统计）；
\item 电子间的 Coulomb 排斥。
\end{enumerate}

作为入门例子，我们研究两个被外势束缚在空间局域位置、并通过 Coulomb 相互作用彼此排斥的电子。我们将发现，两个电子自旋之间的耦合既可以是铁磁的，也可以是反铁磁的，取决于无相互作用态的性质。这些正是相互作用电子体系得以产生各种磁结构的基本机制。

\section{铁磁交换}

考虑两个电子，占据两个近简并的单粒子轨道 $\phi_{1}$、$\phi_{2}$，能量分别为 $\epsilon_{1}$、$\epsilon_{2}$。我们假设谱的其余部分与这两个态之间隔着一个大能隙。图~\ref{fig:2.1} 画出了这类态的一维例子。实际上，它们可以是一个原子或分子的低角动量多重态的成员。我们还假设有两个电子，可以占据其中任一个或两个态。不同组态的简并被 Coulomb 相互作用解除。我们把电子场算符在 $\phi_{i}$ 基组中展开为

\begin{equation*}
\psi_ {s} ^ {\dagger} (x) = \sum_ {i = 1} ^ {2} \phi_ {i} ^ {*} c _ {is} ^ {\dagger}, \quad s = \uparrow , \downarrow
\tag{2.1}
\end{equation*}

\begin{figure}[H]
\centering
\includegraphics[width=0.62\textwidth]{fig2_1_scan.png}
\caption{彼此铁磁耦合的两个电子的态。}
\label{fig:2.1}
\end{figure}

二次量子化形式下的两体 Coulomb 相互作用由（见 (1.17)）

\begin{equation*}
\mathcal {V} = \frac {1}{2} \int d ^ {3} x \, d ^ {3} y \, \tilde {v} (\mathbf {x}, \mathbf {y}) \sum_ {ss'} \psi_ {s} ^ {\dagger} (\mathbf {x}) \psi_ {s '} ^ {\dagger} (\mathbf {y}) \psi_ {s '} (\mathbf {y}) \psi_ {s} (\mathbf {x})
\tag{2.2}
\end{equation*}

给出。用 (2.1) 把 (2.2) 写到 $i$ 表象中：

\begin{equation*}
\mathcal {V} = \sum_ {i \neq i'} U _ {ii'} n _ {i} n _ {i'} + \sum_ {i} U _ {ii} \rho_ {i \uparrow} \rho_ {i \downarrow} + \sum_ {ss',\, i \neq i'} J ^ {F} c _ {is} ^ {\dagger} c _ {i's'} ^ {\dagger} c _ {is'} c _ {i's},
\tag{2.3}
\end{equation*}

其中 $n_i = \sum_s \rho_{is}$。相互作用参数由\emph{直接}积分给出：

\begin{equation*}
U _ {ii'} = \frac {1}{2} \int d ^ {3} x \, d ^ {3} y \, \tilde {v} (\mathbf {x}, \mathbf {y}) \, | \phi_ {i} (\mathbf {x}) | ^ {2} \, | \phi_ {i'} (\mathbf {y}) | ^ {2}
\tag{2.4}
\end{equation*}

以及\emph{铁磁交换}常数：

\begin{equation*}
J ^ {F} = \frac {1}{2} \int d ^ {3} x \, d ^ {3} y \, \tilde {v} (\mathbf {x}, \mathbf {y}) \, \phi_ {i} (\mathbf {x}) \phi_ {i'} ^ {*} (\mathbf {x}) \phi_ {i'} (\mathbf {y}) \phi_ {i} ^ {*} (\mathbf {y}).
\tag{2.5}
\end{equation*}

容易看出 $U_{ii'}$ 为正，$J^{F}$ 为实。显然，对短程相互作用 $\tilde{v} = \delta(\mathbf{x} - \mathbf{y})$，有 $J^{F} > 0$。对长程 Coulomb 相互作用 $v^{c} = e^{2}/|\mathbf{x} - \mathbf{y}|$，$J^{F}$ 的正性证明如下。对称算符 $v^{c}(\mathbf{x}, \mathbf{y})$ 的本征值为

\begin{equation*}
\int d ^ {3} x \, e ^ {i \mathbf {kx}} \frac {e ^ {2}}{| \mathbf {x} |} = \frac {4 \pi e ^ {2}}{\mathbf {k} ^ {2}} > 0.
\tag{2.6}
\end{equation*}

因此 $v^{c}$ 的一切期望值都为正，特别是在态 $\Phi(\mathbf{x}) = \phi_{i'}^{*}\phi_{i}$ 中的期望值，即 (2.5) 为正。这样我们证明了 (2.5) 在两种极限情形下均为正：(i) 完全屏蔽，(ii) 完全无屏蔽。

现在，若

\begin{equation*}
\epsilon_ {1} + \epsilon_ {2} + U _ {12} < \min \left[ 2 \epsilon_ {1} + U _ {11}, \; 2 \epsilon_ {2} + U _ {22} \right],
\tag{2.7}
\end{equation*}

则基态主要含有占据不同轨道的两个电子，即 $n_{i} \approx 1$，低能态由

\begin{equation*}
\left\{ \left| s _ {1}, s _ {2} \right\rangle \right\} , \quad s _ {i} = \uparrow , \downarrow
\tag{2.8}
\end{equation*}

给出。交换相互作用 $J^{F}$ 在 (2.8) 空间中表现为 Heisenberg 相互作用：

\begin{equation*}
J ^ {F} \sum_ {ss'} c _ {is} ^ {\dagger} c _ {i's'} ^ {\dagger} c _ {is'} c _ {i's} = - 2 J ^ {F} \left( \mathbf {S} _ {i} \cdot \mathbf {S} _ {i'} + \frac {1}{4} n _ {i} n _ {i'} \right).
\tag{2.9}
\end{equation*}

$\mathbf{S}_i^\alpha$（$\alpha = x, y, z$）是自旋 $1/2$ 算符的分量（见 (A.19)）：

\begin{equation*}
\mathbf {S} _ {i} \equiv \frac {1}{2} \sum_ {ss'} c _ {is} ^ {\dagger} \vec {\sigma} _ {ss'} c _ {is'}.
\tag{2.10}
\end{equation*}

交换积分 (2.5) 依赖轨道间的空间交叠。(2.3) 与 (2.9) 显式地展示了占据这类轨道的电子自旋之间的铁磁耦合。

一个物理论证：各自旋彼此平行排列、构成对称自旋态后，会削弱 Coulomb 排斥的效应。此时两电子态是反对称的，因而其轨道波函数必须在 $\mathbf{x} - \mathbf{y} = 0$ 处为零——那里正是 Coulomb 势最大的地方。

这一效应出现在微扰论的一阶，与原子中 Hund 定则背后的物理相同：开壳层中的电子占据轨道时使总自旋最大化\footnote{第二 Hund 定则说，第一定则中的含糊之处通过使总轨道角动量最大化来解决。}。

\section{反铁磁交换}

展示两个电子间反铁磁耦合的最简单体系是 $H_{2}$ 分子，最早由 Heitler 和 London 讨论\footnote{这里的讨论相当紧随 Mattis 的书。}。

把两个质子冻结在平衡位置 $R_{i}$，记电子坐标为 $x_{i}$（$i = 1, 2$）。完整哈密顿量为

\begin{equation*}
H = \sum_ {i = 1} ^ {2} H _ {i} ^ {\infty} + \Delta H,
\tag{2.11}
\end{equation*}

\begin{figure}[H]
\centering
\includegraphics[width=0.55\textwidth]{fig2_2.jpg}
\caption{彼此反铁磁耦合的两个电子的态。}
\label{fig:2.2}
\end{figure}

\begin{align*}
H _ {i} ^ {\infty} & = - \frac {\hbar^ {2}}{2 m} \nabla_ {i} ^ {2} - \frac {e ^ {2}}{| \mathbf {x} _ {i} - \mathbf {R} _ {i} |},\\
\Delta H & = - \sum_ {i \neq j} \frac {e ^ {2}}{| \mathbf {x} _ {i} - \mathbf {R} _ {j} |} + \frac {e ^ {2}}{| \mathbf {x} _ {1} - \mathbf {x} _ {2} |} + \frac {e ^ {2}}{| \mathbf {R} _ {1} - \mathbf {R} _ {2} |}.
\tag{2.12}
\end{align*}

两个原子波函数 $\phi_{1}$、$\phi_{2}$ 分别以原子 1、2 为中心，如图~\ref{fig:2.2} 所示。它们满足原子 Schrödinger 方程：

\begin{equation*}
H _ {i} ^ {\infty} \phi_ {i} (\mathbf {x} _ {i}) = e _ {0} \phi_ {i} (\mathbf {x} _ {i}), \quad i = 1, 2.
\tag{2.13}
\end{equation*}

当原子间距有限，$|R_{1}-R_{2}|<\infty$ 时，$\phi_{1}$ 与 $\phi_{2}$ 不正交，其交叠为

\begin{equation*}
\lambda = \int d ^ {3} x \, \phi_ {1} ^ {*} (\mathbf {x}) \phi_ {2} (\mathbf {x}).
\tag{2.14}
\end{equation*}

用 $\phi_{i}$ 构造轨道波函数

\begin{align*}
\psi^ {a} & = \phi_ {1} (\mathbf {x} _ {1}) \phi_ {2} (\mathbf {x} _ {2}),\\
\psi^ {b} & = \phi_ {1} (\mathbf {x} _ {2}) \phi_ {2} (\mathbf {x} _ {1}).
\tag{2.15}
\end{align*}

(2.15) 描述处于不同格点上的电子。我们可用它们构造变分波函数 $\Psi$，以避免大的在位 Coulomb 排斥（稍后我们将用波函数的自旋部分 $\chi$ 施加反对称约束）：

\begin{equation*}
\Psi = \left( c ^ {a} \psi^ {a} + c ^ {b} \psi^ {b} \right) \chi.
\tag{2.16}
\end{equation*}

变分能量为

\begin{equation*}
E = \langle \Psi | H | \Psi \rangle / \langle \Psi | \Psi \rangle.
\tag{2.17}
\end{equation*}

对 $c^a, c^b$ 极小化 $E[c^a, c^b]$，得到矩阵方程：

\begin{equation*}
\left[ \left( \begin{array}{cc} U ^ {d} & U ^ {o} \\ (U ^ {o}) ^ {*} & U ^ {d} \end{array} \right) - (E - 2 e _ {0}) \left( \begin{array}{cc} 1 & \lambda^ {2} \\ \lambda^ {*2} & 1 \end{array} \right) \right] \binom{c ^ {a}}{c ^ {b}} = 0,
\tag{2.18}
\end{equation*}

其中

\begin{equation*}
U ^ {d} = \int d ^ {3} x _ {1} \, d ^ {3} x _ {2} \; \Delta H \, | \psi^ {a} | ^ {2},
\tag{2.19}
\end{equation*}

\begin{equation*}
U ^ {o} = \int d ^ {3} x _ {1} \, d ^ {3} x _ {2} \; \Delta H \, (\psi^ {a}) ^ {*} \psi^ {b}
\tag{2.20}
\end{equation*}

是对角与非对角项。由于 $\psi^{a}$、$\psi^{b}$ 可以取为实（它们由氢原子基态构成），$\lambda$ 与 $U^{o}$ 也可取为实。(2.18) 的解为

\begin{equation*}
E ^ {\pm} = 2 e _ {0} + \frac {U ^ {d} \pm U ^ {o}}{1 \pm \lambda^ {2}},
\tag{2.21}
\end{equation*}

\begin{equation*}
\binom{c ^ {a}}{c ^ {b}} ^ {\pm} = \frac {1}{\sqrt {2}} \binom{1}{\mp 1}.
\tag{2.22}
\end{equation*}

现在施加反对称条件：$\Psi$ 与波函数自旋部分 $\chi(\mathbf{x}_{1},\mathbf{x}_{2})$ 的乘积必须对交换电子坐标反对称：

\begin{align*}
\Psi^ {\pm} & = \frac {1}{\sqrt {2}} \left( \psi^ {a} \mp \psi^ {b} \right) \chi^ {\pm}\\
& = \frac {1}{\sqrt {2}} \left[ \phi_ {1} (\mathbf {x} _ {1}) \phi_ {2} (\mathbf {x} _ {2}) \mp \phi_ {1} (\mathbf {x} _ {2}) \phi_ {2} (\mathbf {x} _ {1}) \right] \chi^ {\pm},\\
\chi^ {\pm} & = \frac {1}{\sqrt {2}} \left[ \chi^ {\uparrow} (\mathbf {x} _ {1}) \chi^ {\downarrow} (\mathbf {x} _ {2}) \pm \chi^ {\uparrow} (\mathbf {x} _ {2}) \chi^ {\downarrow} (\mathbf {x} _ {1}) \right],
\tag{2.23}
\end{align*}

其中 $\chi^{\pm}$ 描述反平行自旋。容易验证 $\Psi^{\pm}$ 是两个 Slater 行列式的线性组合，每个行列式都可写成一个 Fock 态\footnote{严格地说，为了把它们定义为 Fock 态，必须先把 $\phi_{1}$、$\phi_{2}$ 正交化。}。于是，用二次量子化记号可以写

\begin{equation*}
\Psi^ {\pm} = \frac {1}{\sqrt {2}} \left( c _ {1 \uparrow} ^ {\dagger} c _ {2 \downarrow} ^ {\dagger} \pm c _ {1 \downarrow} ^ {\dagger} c _ {2 \uparrow} ^ {\dagger} \right) | 0 \rangle .
\tag{2.24}
\end{equation*}

由这一形式显而易见，$\Psi^{-}$ 与 $\Psi^{+}$ 正是总自旋算符

\begin{equation*}
\mathbf {S} ^ {\mathrm{tot}} = \sum_ {i = 1} ^ {2} \frac {1}{2} \sum_ {ss'} c _ {is} ^ {\dagger} \vec {\sigma} _ {ss'} c _ {is'}
\tag{2.25}
\end{equation*}

的三重态和单态本征态。单态与三重态的劈裂为

\begin{align*}
E ^ {+} - E ^ {-} & = J\\
J & = 2 \, \frac {U ^ {d} \lambda^ {2} - U ^ {o}}{1 - \lambda^ {4}}.
\tag{2.26}
\end{align*}

通过显式计算参数 $U^{d}$ 与 $U^{o}$ 可以证明，对所有 $R_{12}$，$J$ 为正（反铁磁）。因此基态是单态，较高态是三重态，有效哈密顿量可以用 Heisenberg 反铁磁体表示：

\begin{align*}
\mathcal {H} ^ {\mathrm{QHM}} & = + J \, \mathbf {S} _ {1} \cdot \mathbf {S} _ {2}\\
& = J \left[ \frac {1}{2} (\mathbf {S} ^ {\mathrm{tot}}) ^ {2} - \frac {3}{4} \right] = \left\{ \begin{array}{ll} \frac {1}{4} J & S^{\mathrm{tot}} = 1 \\[2pt] - \frac {3}{4} J & S^{\mathrm{tot}} = 0 \end{array} \right. .
\tag{2.27}
\end{align*}

物理论证：反平行的自旋可以利用杂化，在另一个电子占据该格点的同时跳跃过去，从而降低动能。平行自旋则被 Pauli 原理禁止进行这种虚过程，这使三重态能量高于单态。

小结：当两个电子局域在能量相近的轨道上并有排斥相互作用时，它们的磁耦合为：

\begin{enumerate}[itemsep=1pt]
\item \textbf{铁磁}：若两轨道正交但占据空间的同一区域。此时自旋平行排列降低相互作用能。
\item \textbf{反铁磁}：若两轨道不正交但在空间上分离。此时自旋反平行排列降低动能。
\end{enumerate}

\section{超交换}

上一节我们看到，近邻氢原子上的两个电子由于在位排斥与离域能量的相互抗争而倾向于反铁磁耦合。用二次量子化算符，把能量按跳跃矩阵元展开到二阶，可以简单地导出这一效应。考虑局域在两个原子（标记 $i = 1, 2$）上的两个正交轨道。两态之间的隧穿由跳跃哈密顿量描述：

\begin{equation*}
\mathcal {H} ^ {t} = - t \sum_ {s} \left( c _ {1s} ^ {\dagger} c _ {2s} + c _ {2s} ^ {\dagger} c _ {1s} \right).
\tag{2.28}
\end{equation*}

为简单起见，取在位（Hubbard）相互作用：

\begin{equation*}
\mathcal {U} = U \sum_ {i} n _ {i \uparrow} n _ {i \downarrow}.
\tag{2.29}
\end{equation*}

对大的 $U/t$，取 $\mathcal{U}$ 为零级哈密顿量，其基态流形是四重简并的：

\begin{equation*}
\{0 \} = \left| s _ {1}, s _ {2} \right\rangle , \quad s _ {i} = \uparrow , \downarrow, \; i = 1, 2.
\tag{2.30}
\end{equation*}

双占据轨道的能量比多重态 (2.30) 高 $U$。$\mathcal{H}^t$ 的一阶微扰会把我们带出基态流形，因为它把一个电子放到另一个上面。在子空间 (2.30) 内做二阶微扰，得到由标准表达式\footnote{参看 Landau 与 Lifshitz，见本章文献注释。}给出的有效哈密顿量

\begin{equation*}
\langle a | \mathcal {H} ^ {(2)} | b \rangle = - \left\langle a \left| \mathcal {H} ^ {t} \frac {1 - P _ {0}}{\mathcal {U}} \mathcal {H} ^ {t} \right| b \right\rangle = - \sum_ {n \not\in \{0 \}} \langle a | \mathcal {H} ^ {t} | n \rangle \, \frac {1}{\langle n | \mathcal {U} | n \rangle} \, \langle n | \mathcal {H} ^ {t} | b \rangle ,
\tag{2.31}
\end{equation*}

其中 $a, b$ 是子空间 $\{0\}$ 中的态，$P_{0}$ 是其投影算符，而 $|n\rangle$ 含有双占据。和式中的每一项都可以用一条“交换路径”表示。例如

\begin{equation*}
\begin{array}{ccccc}
 & \mathcal {H} ^ {t} & \ \ \frac {1}{\mathcal {U}}\ \ & & \mathcal {H} ^ {t} \\
| \uparrow , \downarrow \rangle & \to & | \uparrow \downarrow , 0 \rangle & \to & | \downarrow , \uparrow \rangle \\
 & & \to & | 0, \uparrow \downarrow \rangle & \to\ | \downarrow , \uparrow \rangle .
\end{array}
\tag{2.32}
\end{equation*}

(2.32) 中两条路径各给出贡献 $-t^{2}/U$。但也存在被 Pauli 不相容原理阻塞的路径，例如

\begin{equation*}
\begin{array}{ccc}
 & \mathcal {H} ^ {t} & \\
| \uparrow , \uparrow \rangle & \to & 0.
\end{array}
\tag{2.33}
\end{equation*}

因此三重态不能靠虚双占获得二阶能量增益。连接这些初末态的算符可以写成 $S=\tfrac{1}{2}$ 自旋算符的乘积 $S_{1}^{-}S_{2}^{+}$。把描述 $S_{2}^{-}S_{1}^{+}$ 与 $S_{1}^{z}S_{2}^{z}$ 矩阵元的路径也包括进来，我们发现基态子空间中的有效哈密顿量可写成各向同性的反铁磁 Heisenberg 交换：

\begin{equation*}
\mathcal {H} ^ {(2)} = J \, \mathbf {S} _ {1} \cdot \mathbf {S} _ {2} \, , \quad J = \frac {4 t ^ {2}}{U}.
\tag{2.34}
\end{equation*}

这一讨论容易推广到多格点体系：任何两个由跳跃项 $H^{t}$ 耦合的自旋之间都会产生反铁磁耦合。

虚双占据过程称为“\textbf{超交换}”（superexchange）。上一节对氢分子的分析用变分方法描述了同一效应。我们要把超交换与 2.1 节讲的铁磁直接交换或 Hund 定则区分开：后者是相互作用一阶微扰的结果\footnote{Anderson 在其原始论文（见文献注释）中考虑了几种交换机制。在那里，“超交换”专指经过一个中间离子的耦合。}。

\begin{figure}[H]
\centering
\includegraphics[width=0.5\textwidth]{fig2_3a.jpg}\\[4pt]
\includegraphics[width=0.5\textwidth]{fig2_3b.jpg}
\caption{铜氧化物中的超交换。(a)~低能流形的一个成员；(b)~超交换路径中的一个高能组态。}
\label{fig:2.3}
\end{figure}

超交换耦合在许多其他带有未配对自旋的体系中产生。当前备受关注的著名例子是铜氧化物反铁磁体。未配对自旋位于 $Cu^{2+}$ 离子上，二者之间有一个氧“桥”离子，平均电离态接近 $O^{2-}$。这样的三原子组合见图~\ref{fig:2.3}。

若铜与氧之间没有跳跃，则键上的简并基态流形为

\begin{equation*}
\{ | 0 \rangle \} = \left\{ \left| \mathrm{Cu} ^ {2+} s _ {1}, \mathrm{O} ^ {2-}, \mathrm{Cu} ^ {2+} s _ {2} \right\rangle , \quad s _ {i} = \uparrow , \downarrow \right\}
\tag{2.35}
\end{equation*}

解除 (2.35) 简并的基态能量的最低阶修正是铜--氧跳跃的四阶项。超交换路径的中间态及其能量的例子有

\begin{equation*}
\left| \mathrm{Cu} ^ {1+}, \mathrm{O} ^ {2-}, \mathrm{Cu} ^ {3+} \uparrow \downarrow \right\rangle , \quad E = U _ {d},
\tag{2.36}
\end{equation*}

以及另一个具有不同能量分母的态

\begin{equation*}
\left| \mathrm{Cu} ^ {1+}, \mathrm{O} \uparrow \downarrow , \mathrm{Cu} ^ {1+} \right\rangle , \quad E = 2 (\epsilon_ {p} - \epsilon_ {d}) + U _ {p},
\tag{2.37}
\end{equation*}

其中 $\epsilon_{d(p)}$ 与 $U_{d(p)}$ 分别是铜（氧）上的单粒子能量和相互作用能。如 (2.33) 所示，平行自旋态被阻挡，无法跳跃进入这两个中间态中的任何一个。这就解除了 (2.35) 中单态与三重态的简并，简并的解除由 Heisenberg 哈密顿量描述。反铁磁耦合由对所有超交换路径求和给出，因此由能量最低的中间态主导。

\begin{exercises}

\item 自旋算符的费米子表示由 (2.10) 给出。证明重要恒等式 (2.9)：

\begin{equation*}
\sum_ {ss'} c _ {1s} ^ {\dagger} c _ {2s'} ^ {\dagger} c _ {1s'} c _ {2s} = - \frac {1}{2} \left\{ n _ {1} n _ {2} + 4 \, \mathbf {S} _ {1} \cdot \mathbf {S} _ {2} \right\},
\tag{2.38}
\end{equation*}

其中 $n_i = \sum_s c_{is}^\dagger c_{is}$，$i = 1,2$。

\item 证明 $\mathbf{S}_i$ 满足角动量对易关系。求 $\mathbf{S}_i^2$ 作为 $n_i$ 函数的简单表达式。

\end{exercises}

\begin{chapbib}
\item 氢分子的分析相当紧随 Mattis 的书：
  D. C. Mattis, \textit{The Theory of Magnetism I} (Springer-Verlag, 1988)。
\item 超交换理论由 Kramers 提出，并由 Anderson 在其开创性论文中发展：
  P. W. Anderson, Phys. Rev. \textbf{79}, 350 (1950)。
\item 铜氧化物超导体的超交换耦合的计算，例如：
  F. C. Zhang and T. M. Rice, Phys. Rev. B \textbf{37}, 3759 (1988)。
\item 微扰论基础参看：
  L. D. Landau and E. M. Lifshitz, \textit{Quantum Mechanics} (Pergamon Press, 1977), 第六章。
\end{chapbib}
